%% ijbm_SI.tex -- Supplementary material (Appendix A) for ijbm_main.tex,
%% International Journal of Biological Macromolecules format. Text copied unchanged from manuscript_short/grip_finding_SI.tex.
\documentclass[preprint,12pt]{elsarticle}
\usepackage{amsmath,amssymb,booktabs,graphicx}
\usepackage[hidelinks]{hyperref}
\graphicspath{{figures/}}
\renewcommand{\thesection}{S\arabic{section}}
\renewcommand{\thetable}{S\arabic{table}}
\renewcommand{\theequation}{S\arabic{equation}}
\renewcommand{\thefigure}{S\arabic{figure}}
\journal{International Journal of Biological Macromolecules}
\begin{document}
\begin{frontmatter}
\title{Supplementary material for: Fixed and variable sugar contacts at the subsites of
starch-degrading enzymes: a structural basis for processing functionalised starch}
\author[a,b]{Aldo F. Combariza}
\ead{aldo.combariza@unisucre.edu.co}
\affiliation[a]{organization={Grupo de Investigaci\'{o}n {\itshape in silico},
  Departamento de Biolog\'{\i}a y Qu\'{\i}mica, Universidad de Sucre},
  addressline={Carrera 28 No.\ 5-267},
  city={Sincelejo},
  state={Sucre},
  country={Colombia}}
\affiliation[b]{organization={Doctorado en Ciencias F\'{\i}sicas, Sistema Universitario Estatal del Caribe (SUE Caribe)},
  country={Colombia}}
\end{frontmatter}


\section{Ring distortion at subsite $-1$}
The ring at subsite $-1$ is bent away from the chair more often than anywhere else. With
the Cremer--Pople polar angle $\theta$ folded to its distance from the nearest
chair~\citep{Cremer1975}, and a ring called distorted above 20$^\circ$, 8 of 23
density-supported rings at $-1$ are distorted in GH13 (8 entries from 6 of 14 proteins:
1QHO, 1QHP, 1QI5, 3BCD, 4J3V, 5TD4, 6JFJ, 6JFX), against 5 of 97 at the other subsites;
in GH14, 4 of 15 at $-1$ (both proteins) against 0 of 21; in GH57, none of 14. Distortion of
the $-1$ ring toward the transition-state geometry is expected for
glycosidases~\citep{Davies2012}; here it is found where the fixed hydrogen-bonding
residues are. Two of the eight GH13 entries, 1QI5 and 5TD4, are engineered Asp$\to$Asn
mutants of the residue at the Asp289 position (main text). At the other subsites, GH31 has
5 distorted rings of 38 and GH57 2 of 38 (GH31 has no density-supported ring at $-1$).
Source: \texttt{42\_ring\_distortion/summary.tsv}.

\section{Conserved water sites at the catalytic residues}
Conserved water sites sit at the catalytic carboxylates in four
families (water sites occupied in at least half of the superposed entries, within
3.6~\AA{} of a catalytic oxygen): in GH13, W3 at 3.09~\AA{} from the nucleophile Asp179
(251 of 411 entries) and W4 at 2.43~\AA{} from the acid/base Glu204 (235 of 411); in GH14,
W34 at 3.00~\AA{} from the base Glu380 (35 of 50) and W123 at 3.10~\AA{} from the acid
Glu186 (26 of 50); in GH31, W17 at 2.84~\AA{} from the nucleophile Asp443 (50 of 73) and W1
and W12 at 2.85 and 2.33~\AA{} from the acid Asp542 (58 and 52 of 73); in GH57, W9 at
2.86~\AA{} from the nucleophile Glu123 (20 of 25) and W170 at 3.09~\AA{} from the acid Asp214
(14 of 25). In GH13 the fraction depends on resolution (Table~\ref{tab:sires}): W3 is occupied in
75.5\% and W4 in 83.0\% of entries better than 1.6~\AA, and in 18.6\% and 12.9\% of entries
at 2.6~\AA{} or worse; the fall is not monotonic for W3 (62.7\% at 1.6--1.9~\AA, 80.2\% at
1.9--2.2~\AA). The pooled fractions above are therefore below the occupancy seen in the best
data. The hydrogen-bond network that the census finds
fixed at the catalytic subsite thus extends from the protein side chains to waters bound to
the catalytic residues. Whether these waters take part in the reaction is not decided by
these data.

\begin{table}[h]
\centering\small
\caption{Occupancy of the GH13 water sites W3 and W4 by resolution band of the entry.
Source: \texttt{resultados/analysis/conserved\_sites\_resolution\_confound.md}.}
\label{tab:sires}
\begin{tabular}{lrrr}
\toprule
Resolution & Entries & W3 occupied & W4 occupied\\
\midrule
$<1.6$~\AA & 53 & 75.5\% & 83.0\%\\
1.6--1.9~\AA & 75 & 62.7\% & 81.3\%\\
1.9--2.2~\AA & 131 & 80.2\% & 62.6\%\\
2.2--2.6~\AA & 83 & 55.4\% & 47.0\%\\
$\ge2.6$~\AA & 70 & 18.6\% & 12.9\%\\
\bottomrule
\end{tabular}
\end{table}

\section{Computed pK$_\mathrm{a}$ of the catalytic residues}
The computed protonation state of the catalytic residues was compared between chains with
and without a bound molecule near them (PROPKA; the model pK$_\mathrm{a}$ of a free Asp/Glu in water is
3.80/4.50, and the value reported is that estimate adjusted for the surrounding
structure). In GH13, the acid/base residue's median computed pK$_\mathrm{a}$ is 6.93 in
chains with no sugar or other molecule within 4.0~\AA{} of it and 5.30 in chains with one
(314 and 196 chains); the nucleophile shifts the other way, 6.63 without and 7.46 with (308
and 203 chains). All values are in Table~\ref{tab:propka}. The acid of GH31 shifts as in GH13 (7.76
without, 5.78 with); the GH31 nucleophile does not shift (5.11 and 5.16); in GH57 both
residues are lower with a ligand (acid 6.20 and 5.39, nucleophile 8.35 and 7.41). The
chains with and without a ligand are different crystal forms and in part different proteins,
so the comparison is not paired. These are computed estimates (PROPKA), not measurements,
and they do not give a mechanism; the rise of the GH13 nucleophile with a ligand is the
opposite of the simplest expectation and is not explained here.

\begin{table}[h]
\centering\small
\caption{Median computed pK$_\mathrm{a}$ (PROPKA) of the catalytic residues in chains with
and without a sugar or other molecule within 4.0~\AA{} of the residue; number of chains in
parentheses. Source: \texttt{resultados/calc/01\_propka/propka\_all.tsv}.}
\label{tab:propka}
\begin{tabular}{llll}
\toprule
Family & Residue & Without ligand & With ligand\\
\midrule
GH13 & nucleophile Asp & 6.63 (308) & 7.46 (203)\\
GH13 & acid/base Glu & 6.93 (314) & 5.30 (196)\\
GH31 & nucleophile Asp & 5.11 (50) & 5.16 (81)\\
GH31 & acid Asp & 7.76 (33) & 5.78 (88)\\
GH57 & nucleophile Glu & 8.35 (19) & 7.41 (35)\\
GH57 & acid Asp & 6.20 (19) & 5.39 (27)\\
\bottomrule
\end{tabular}
\end{table}

\section{MM/GBSA of four inhibitor complexes of human lysosomal $\alpha$-glucosidase}
The kinetic values of the main text are measured for whole hydroxyls, not for single
residues. A per-residue split can be computed for the
one GH31 enzyme with several deposited inhibitor complexes and measured potencies: human
lysosomal $\alpha$-glucosidase (GAA, UniProt P10253), with 1-deoxynojirimycin (5NN5),
miglitol (5NN6), a 2-pentyl deoxynojirimycin analogue (9GTC) and a 2-heptyl pyrrolidine
(9GTW). Each complex (chain A, ring nitrogen of the inhibitor protonated) was built with
ff14SB and GAFF2/AM1-BCC, minimised in implicit solvent (sander, AmberTools~26; 1000 steps, 300 of steepest descent;
igb~5, 16~\AA{} cut-off) with protein heavy atoms beyond 6~\AA{} of the inhibitor restrained
at 5~kcal\,mol$^{-1}$\,\AA$^{-2}$, and evaluated once with MMPBSA.py (generalised Born
model igb~5, 0.15~M salt):
\begin{equation}
\Delta G_\mathrm{bind} = G_\mathrm{complex}-G_\mathrm{receptor}-G_\mathrm{ligand}
= \Delta E_\mathrm{vdW}+\Delta E_\mathrm{elec}+\Delta G_\mathrm{GB}+\Delta G_\mathrm{SA},
\label{eq:gbsa}
\end{equation}
where $\Delta E_\mathrm{vdW}$ and $\Delta E_\mathrm{elec}$ are the van der Waals and
Coulomb energies between inhibitor and protein, $\Delta G_\mathrm{GB}$ is the change in
polar solvation energy, and $\Delta G_\mathrm{SA}$ the non-polar term from the buried surface;
no entropy term is included. For 5NN5: $-16.49-427.80+384.40-3.91=-63.80$~kcal\,mol$^{-1}$.
The per-residue decomposition gives half of each inhibitor--residue pair energy to the residue;
the residue contributions at nine GH31 positions of Table~1 of the main text (Asp327, His600,
Asp443, Asp542, Trp441, Trp539, Phe575, Ile328 and Ile364; 2QMJ numbering), mapped to
the 2QMJ reference through the contact census of 5NN5, are summed by kind of contact in
Table~\ref{tab:gbsa}. The ``protein'' column is the sum over the 26--30 residues within
6~\AA{} of the inhibitor that were decomposed, not over the whole protein, so inhibitor and
protein columns do not add up to $\Delta G_\mathrm{bind}$. In all four complexes the four
hydrogen-bonding positions carry $-11.9$ to $-15.9$~kcal\,mol$^{-1}$ ($-3.0$ to
$-4.0$ per position), 46--58\% of the decomposed protein residues; the
three aromatic positions $-2.5$ to $-3.2$ ($-0.8$ to $-1.1$ per position); the two
aliphatic positions $-1.1$ to $-1.6$ ($-0.55$ to $-0.79$ per position). The
largest single contribution is Asp404 (Asp327 of the reference), $-5.9$ to
$-7.3$~kcal\,mol$^{-1}$. The one exception within a group is His674 (His600), which gives
$-2.7$ to $-3.3$~kcal\,mol$^{-1}$ to the three six-membered iminosugars and $-0.05$ to the
five-membered pyrrolidine of 9GTW. The total $\Delta G_\mathrm{bind}$ does not follow the
measured potency: $-63.8$ to $-67.7$~kcal\,mol$^{-1}$ against pChEMBL 5.68--7.72
(Pearson $r = +0.07$ between pChEMBL and $\Delta G_\mathrm{bind}$, Spearman $\rho=0.0$,
$n=4$; the pChEMBL values mix $K_\mathrm{i}$ and IC$_{50}$ from different assays, and three of
the four are single measurements). The limits are stated flatly: one minimised structure
per complex, no sampling and no entropy term, so the energies are not binding free
energies; the inhibitors carry a protonated ring nitrogen, so the large terms of the three
Asp positions are mainly charge--charge interactions with the ammonium, and the order
H $>$ Ar $>$ Al would follow from the charges for almost any cationic ligand in this pocket.
The calculation therefore does not test the contact-type rule.

\begin{table}[h]
\centering\small
\caption{MM/GBSA of four GAA inhibitor complexes (kcal\,mol$^{-1}$): total
$\Delta G_\mathrm{bind}$ (eq~\ref{eq:gbsa}), the inhibitor and protein halves of the
per-residue decomposition, and the sum over the nine GH31 positions of Table~1 of each kind of
contact (H: Asp327, His600, Asp443, Asp542; Ar: Trp441, Trp539, Phe575; Al: Ile328, Ile364;
2QMJ numbering; GAA residues Asp404, His674, Asp518, Asp616, Trp516, Trp613, Phe649,
Leu405, Ile441). pChEMBL: median against GAA. Sources:
\texttt{resultados/calc/09\_mmgbsa/series/}, \texttt{resultados/analysis/43\_mmgbsa\_subsite/}.}
\label{tab:gbsa}
\begin{tabular}{llrrrrrrr}
\toprule
PDB & Inhibitor & pChEMBL & $\Delta G_\mathrm{bind}$ & Inhibitor & Protein & H & Ar & Al\\
\midrule
5NN5 & 1-deoxynojirimycin & 6.82 & $-63.80$ & $-34.41$ & $-27.18$ & $-15.89$ & $-2.71$ & $-1.10$\\
5NN6 & miglitol & 6.46 & $-67.66$ & $-38.12$ & $-26.47$ & $-12.89$ & $-2.53$ & $-1.57$\\
9GTC & 2-pentyl-DNJ analogue & 7.72 & $-65.73$ & $-38.29$ & $-25.82$ & $-14.66$ & $-2.86$ & $-1.30$\\
9GTW & 2-heptyl pyrrolidine & 5.68 & $-65.32$ & $-38.05$ & $-25.73$ & $-11.91$ & $-3.16$ & $-1.49$\\
\bottomrule
\end{tabular}
\end{table}


\bibliographystyle{elsarticle-num}
\bibliography{ijbm}
\end{document}
